\documentclass[11pt]{article}

\usepackage{math_article}
\usepackage{series_ra}

\renewcommand{\ArticleNumeral}{}
\renewcommand{\ArticleTitle}{Uniform Integrability and the\\Vitali Convergence Theorem}
\renewcommand{\ArticleAuthor}{Anqiao Ouyang}
\renewcommand{\ArticleDate}{August 2026}

\begin{document}

\maketitle


\section*{Introduction}


When studying the various modes of convergence for sequences of functions, it is natural to harbor an expectation:

if $f_n$ approaches $f$ in ever stronger senses, shouldn't the integrals approach each other as well?

In other words, when can we pass from

\[
f_n\to f
\]

to

\[
\int_X f_n\,d\mu
\longrightarrow
\int_X f\,d\mu?
\]

This question is a familiar one in Lebesgue integration theory. The monotone convergence theorem settles it by means of nonnegativity and monotonicity; the dominated convergence theorem instead requires a single integrable dominating function $g$ with

\[
|f_n|\leq g.
\]

All of these conditions do the same job: they prevent the ``mass'' of the sequence from escaping in some way.

But the requirement of a dominating function is sometimes too strong.

A sequence of functions may admit no common integrable upper bound at all and still have perfectly good integral behavior. What really matters then is not any single dominating function, but a uniform control of the whole family over the ``regions where the function values are large.''

This is \textbf{uniform integrability}.


\section{Concentrated Mass}


Let us start with the simplest example.

On $([0,1],\mathcal B,m)$ define

\[
f_n(x)
=
n\mathbf 1_{(0,1/n)}(x).
\]

For any fixed $x>0$, once $n$ is large enough,

\[
x\notin(0,1/n),
\]

and hence

\[
f_n(x)\to0.
\]

Therefore

\[
f_n\to0
\qquad\text{a.e.}
\]

At the same time, for any $\varepsilon>0$ and $n>\varepsilon$,

\[
\{|f_n|>\varepsilon\}
=
(0,1/n),
\]

so that

\[
m(\{|f_n|>\varepsilon\})
=
\frac1n\to0.
\]

Hence the sequence also converges to $0$ in measure.

However,

\[
\int_0^1|f_n|\,dx = n\cdot\frac{1}{n} = 1.
\]

Thus

\[
\|f_n\|_{L^1}=1
\]

stays constant, and a fortiori we cannot have

\[
f_n\to0
\quad\text{in }L^1.
\]

What goes wrong?

It is not that the total integral of the sequence grows. In fact it is always equal to $1$.

What actually changes is that this integral is squeezed onto smaller and smaller sets:

\[
\text{height }n
\qquad
\times
\qquad
\text{width }\frac1n.
\]

The sets shrink while the function values grow, and the two effects exactly cancel.

Consequently, merely knowing

\[
\sup_n\|f_n\|_1<\infty
\]

is not enough.

We need to rule out this phenomenon of ``a finite amount of mass concentrating in infinitely tall spikes.''


\section{Definition}


Let $(X,\mathcal F,\mu)$ be a measure space and let $\mathcal G\subset L^1(\mu)$ be a family of integrable functions.

If

\[
\lim_{K\to\infty}
\sup_{f\in\mathcal G}
\int_{\{|f|>K\}}|f|\,d\mu
=
0,
\]

then $\mathcal G$ is said to be \textbf{uniformly integrable}.

For a sequence $\{f_n\}$ this reads

\[
\lim_{K\to\infty}
\sup_n
\int_{\{|f_n|>K\}}
|f_n|\,d\mu
=
0.
\]

The definition is concerned not with the whole integral

\[
\int_X|f_n|\,d\mu,
\]

but with the part where the function values are very large:

\[
\int_{\{|f_n|>K\}}|f_n|\,d\mu.
\]

For a fixed large $K$, each function may well have a part exceeding $K$. But if the family is uniformly integrable, then as soon as $K$ is large enough, the integral mass carried by these parts can be made arbitrarily small \textbf{simultaneously}.

The crucial ingredient here is the

\[
\sup_n.
\]

For any single $f_n\in L^1$ we already have

\[
\int_{\{|f_n|>K\}}|f_n|\,d\mu\to0.
\]

Uniform integrability demands that this convergence hold uniformly over the whole family.

It is therefore not just another way of saying ``each function is integrable,'' but a genuine \textbf{property of the family}. In probability theory, uniform integrability is usually defined by exactly the same tail-integral condition.


\section{Boundedness}


First of all, uniform integrability implies

\[
\sup_{f\in\mathcal G}
\|f\|_1<\infty.
\]

The proof is not complicated.

By uniform integrability we can choose some $K>0$ such that

\[
\sup_{f\in\mathcal G}
\int_{\{|f|>K\}}|f|\,d\mu
<1.
\]

Split the integral into two parts:

\[
\int_X|f|\,d\mu
=
\int_{\{|f|\leq K\}}|f|\,d\mu
+
\int_{\{|f|>K\}}|f|\,d\mu.
\]

If $\mu(X)<\infty$, the first term satisfies

\[
\int_{\{|f|\leq K\}}|f|\,d\mu
\leq
K\mu(X).
\]

Hence

\[
\|f\|_1
\leq
K\mu(X)+1.
\]

Therefore

\[
\sup_{f\in\mathcal G}\|f\|_1<\infty.
\]

The converse, however, fails.

The sequence

\[
f_n=n\mathbf1_{(0,1/n)}
\]

from the previous section is already a counterexample. Although

\[
\sup_n\|f_n\|_1=1,
\]

as soon as $n>K$ we have

\[
\{|f_n|>K\}
=
(0,1/n),
\]

and so

\[
\int_{\{|f_n|>K\}}|f_n|\,dx
=
1.
\]

Thus for every $K$,

\[
\sup_n
\int_{\{|f_n|>K\}}|f_n|\,dx
=
1.
\]

The sequence is not uniformly integrable.

So

\[
\boxed{
\text{uniformly integrable}
\Longrightarrow
L^1\text{-bounded}
}
\]

while the converse does not hold.

What separates the two is precisely control over the \textbf{concentration of mass}.


\section{Small Sets}


On a finite measure space, uniform integrability admits another very important interpretation.

If $\mathcal G$ is uniformly integrable, then for every $\varepsilon>0$ there exists $\delta>0$ such that for every measurable set $E$ with

\[
\mu(E)<\delta
\]

we have

\[
\sup_{f\in\mathcal G}
\int_E|f|\,d\mu
<
\varepsilon.
\]

In words:

\begin{quote}
A sufficiently small set cannot carry a large amount of the $L^1$ mass of any function in the family.
\end{quote}


The proof starts directly from truncation.

By uniform integrability, choose $K$ such that

\[
\sup_{f\in\mathcal G}
\int_{\{|f|>K\}}|f|\,d\mu
<
\frac{\varepsilon}{2}.
\]

For any measurable set $E$,

\[
\int_E|f|\,d\mu
=
\int_{E\cap\{|f|\leq K\}}|f|\,d\mu
+
\int_{E\cap\{|f|>K\}}|f|\,d\mu.
\]

The first term is at most

\[
K\mu(E),
\]

and the second is at most

\[
\frac{\varepsilon}{2}.
\]

So it suffices to take

\[
\delta=\frac{\varepsilon}{2K};
\]

whenever $\mu(E)<\delta$ we then get

\[
\int_E|f|\,d\mu<\varepsilon.
\]

This is in fact closer to the geometric meaning of uniform integrability than the tail-integral definition.

Ordinary $L^1$-boundedness only controls

\[
\int_X|f|,
\]

whereas uniform integrability further requires that \textbf{integral mass cannot concentrate arbitrarily on smaller and smaller sets.}

On a finite measure space, this ``uniform absolute continuity on small sets,'' together with uniform $L^1$-boundedness, conversely characterizes uniform integrability.


\section{$L^p$ Control}


Uniform integrability may still look somewhat abstract, but there is a very commonly used sufficient condition.

Let

\[
1<p<\infty
\]

and suppose the family $\mathcal G$ satisfies

\[
\sup_{f\in\mathcal G}\|f\|_p\leq C.
\]

Then $\mathcal G$ is uniformly integrable (in the $L^1$ sense).

Indeed, on the set

\[
A_K=\{|f|>K\}
\]

we have

\[
|f|
\leq
\frac{|f|^p}{K^{p-1}},
\]

and hence

\[
\int_{A_K}|f|\,d\mu
\leq
\frac1{K^{p-1}}
\int_{A_K}|f|^p\,d\mu
\leq
\frac{\|f\|_p^p}{K^{p-1}}.
\]

Therefore

\[
\sup_{f\in\mathcal G}
\int_{\{|f|>K\}}|f|\,d\mu
\leq
\frac{C^p}{K^{p-1}}.
\]

Since $p>1$,

\[
K^{1-p}\to0,
\]

and uniform integrability follows.

This also shows why $p=1$ is a special, critical position.

If all we have is

\[
\sup_n\|f_n\|_1<\infty,
\]

the right-hand side produces no factor that vanishes as $K\to\infty$.

But as soon as there is even a little extra integrability,

\[
p>1,
\]

we obtain the factor

\[
K^{1-p},
\]

which automatically controls the tail of large function values.


\section{Dominated Convergence}


The hypothesis of the dominated convergence theorem also yields uniform integrability automatically.

Suppose there exists

\[
g\in L^1(\mu)
\]

such that for all $n$,

\[
|f_n|\leq g
\qquad\text{a.e.}
\]

If

\[
|f_n|>K,
\]

then necessarily also

\[
g>K.
\]

Hence

\[
\{|f_n|>K\}
\subseteq
\{g>K\},
\]

and

\[
\int_{\{|f_n|>K\}}|f_n|\,d\mu
\leq
\int_{\{g>K\}}g\,d\mu.
\]

Since $g\in L^1$,

\[
\int_{\{g>K\}}g\,d\mu\to0.
\]

Therefore

\[
\{f_n\}
\]

is uniformly integrable.

Thus the common dominating function in the dominated convergence theorem actually provides more structure than is needed to prove convergence of the integrals.

What it gives is

\[
|f_n|\leq g\in L^1,
\]

whereas what is really responsible for controlling the tails of the integrals is

\[
\{f_n\}\text{ uniformly integrable}.
\]

This is one of the reasons why the Vitali convergence theorem is more flexible than the dominated convergence theorem. Uniform integrability can be viewed as a relaxation of the integrable-domination condition.


\section{Vitali's Theorem}


We can now return to the original question.

Suppose

\[
\mu(X)<\infty
\]

and

\[
f_n\to f
\qquad\text{in measure}.
\]

Convergence in measure alone does not imply $L^1$ convergence; the spike functions have already shown this.

The Vitali convergence theorem states that the missing condition is precisely uniform integrability:

\[
\boxed{
f_n\to f\text{ in measure}
+
\{f_n\}\text{ uniformly integrable}
\Longrightarrow
f_n\to f\text{ in }L^1.
}
\]

In fact, within this framework one also obtains a corresponding converse: when $f_n\to f$ in measure, $L^1$ convergence and uniform integrability are essentially equivalent. The version of Vitali's theorem for probability measures is usually stated in exactly this form.

Why does the theorem hold?

Take $\eta>0$ and define the bad set

\[
A_n
=
\{|f_n-f|>\eta\}.
\]

Since

\[
f_n\to f
\qquad\text{in measure},
\]

we have

\[
\mu(A_n)\to0.
\]

Now split the $L^1$ error as

\[
\int_X|f_n-f|\,d\mu
=
\int_{X\setminus A_n}|f_n-f|\,d\mu
+
\int_{A_n}|f_n-f|\,d\mu.
\]

On $X\setminus A_n$,

\[
|f_n-f|\leq\eta,
\]

so

\[
\int_{X\setminus A_n}|f_n-f|\,d\mu
\leq
\eta\mu(X).
\]

Meanwhile the measure of $A_n$ tends to zero.

Uniform integrability guarantees exactly that on such shrinking sets,

\[
\int_{A_n}|f_n|
\]

as well as the corresponding contribution of the limit function cannot retain a fixed positive mass.

Hence the second term tends to zero as well.

Finally, letting first $n\to\infty$ and then $\eta\to0$, we obtain

\[
\|f_n-f\|_1\to0.
\]

This proof reveals the true structure of Vitali's theorem:

\[
\boxed{
\begin{array}{c}
\text{convergence in measure}\\
\Downarrow\\
\text{measure of the bad set}\to 0
\end{array}
\qquad+\qquad
\begin{array}{c}
\text{uniform integrability}\\
\Downarrow\\
\text{small sets cannot carry much integral}
\end{array}
}
\]

The two pieces fit together into exactly $L^1$ convergence.


\section{Two Kinds of Control}


Looking back at the dominated convergence theorem from here, we find that the logic of the two theorems is in fact very close.

Dominated convergence uses one fixed function

\[
g\in L^1
\]

to control all the $f_n$ at once.

This is \textbf{pointwise control}:

\[
|f_n(x)|\leq g(x).
\]

Uniform integrability does not require such a $g$ to exist.

It only requires that the tails of the family be controlled uniformly in the integral sense:

\[
\sup_n
\int_{\{|f_n|>K\}}|f_n|
\to0.
\]

This is \textbf{mass control}.

Consequently there can be uniformly integrable sequences for which no common $L^1$ dominating function exists. Vitali's theorem still handles such cases; this is exactly where it strictly goes beyond the direct scope of the dominated convergence theorem. Stanford's lecture notes on probability theory also give examples of this kind: ``uniformly integrable, but with no common integrable dominating variable.''

So the two theorems can be understood within one logical framework:

\[
\text{DCT}
:
\quad
\text{common }L^1\text{ domination}
\Longrightarrow
\text{uniform integrability}
\Longrightarrow
L^1\text{ limit under control},
\]

and Vitali starts directly from the middle layer.


\section{Infinite Measure}


There is one more boundary here that is easy to overlook.

The finite-measure assumption above is not entirely inessential.

On

\[
(\mathbb R,\mathcal B,m)
\]

consider

\[
f_n(x)
=
\frac1n\mathbf1_{[0,n]}(x).
\]

Clearly

\[
\|f_n\|_1=1.
\]

Moreover, since

\[
\|f_n\|_\infty=\frac1n,
\]

in fact

\[
f_n\to0
\]

even uniformly, so of course also in measure.

On the other hand, for $K>1$,

\[
\{|f_n|>K\}=\varnothing
\]

for all $n$, so by the tail-integral definition $\{f_n\}$ is uniformly integrable.

Yet

\[
\|f_n\|_1=1
\]

still does not tend to zero.

Here the mass does not concentrate on smaller and smaller sets; instead, a different phenomenon occurs:

\[
\text{the mass spreads out over larger and larger regions of the space.}
\]

On a finite measure space there is not enough room for this behavior; on an infinite measure space it requires additional control.

Hence, for a general measure space, controlling only the tail in the direction

\[
|f_n|\to\infty
\]

is not enough; one must also prevent the integral mass from escaping to the ``far away'' parts of the space.

General Vitali-type results therefore add a tightness condition: one can find a region of finite measure outside of which the integral mass is uniformly small for all $n$. The version for general measure spaces genuinely needs this spatial kind of control in addition to control of large values.

The two phenomena are best kept apart:

\[
\boxed{
\begin{aligned}
\text{concentration}
&:\quad
\text{mass squeezes into smaller and smaller sets},\\
\text{escape}
&:\quad
\text{mass runs off to larger and larger regions of space}.
\end{aligned}
}
\]

Uniform integrability rules out the former first; on infinite measure spaces the latter must be dealt with in addition.


\section*{Summary}


Uniform integrability is not about whether a single function is integrable, but about whether the integral mass of a family of functions is uniformly controlled.

It requires

\[
\lim_{K\to\infty}
\sup_n
\int_{\{|f_n|>K\}}
|f_n|\,d\mu
=
0,
\]

that is, a sequence cannot preserve a fixed amount of integral mass through ever taller and ever narrower spikes.

Consequently,

\[
\sup_n\|f_n\|_1<\infty
\]

by itself is not sufficient, whereas uniform $L^p$ control for any $p>1$ automatically yields uniform integrability.

More importantly, on a finite measure space the Vitali convergence theorem ties it to the convergence behavior of the sequence:

\[
\boxed{
f_n\to f\text{ in measure}
\quad+\quad
\{f_n\}\text{ uniformly integrable}
\quad\Longrightarrow\quad
f_n\to f\text{ in }L^1.
}
\]

Convergence in measure tells us that the points where the error is large become fewer and fewer.

Uniform integrability tells us that these ever fewer points cannot still carry a large amount of integral.

The two conditions control two different problems, and together they supply exactly the missing step from convergence in measure to $L^1$ convergence.

From this point of view, uniform integrability is not yet another ``mode of convergence.''

It is closer to a \textbf{compactness condition}: it controls how the mass of a family of functions is distributed, and decides when we may safely pass the limit inside the integral sign.

\end{document}
